\begin{document}
\hypersetup{pdftitle={Mathématiques Générales (Licence 1 & 2) : support de cours},
            pdfauthor={\Auteur},
            pdfsubject={Ensembles, logique et algèbre de Boole, équations et inéquations, systèmes linéaires et programmation linéaire, mesures des angles et trigonométrie, suites numériques},
            pdfkeywords={ensembles, diagramme de Venn, logique, table de vérité, algèbre de Boole, équation du premier degré, équation du second degré, discriminant, inéquation, tableau de signes, système linéaire, Cramer, Gauss, programmation linéaire, régionnement du plan, degré, radian, cercle trigonométrique, sinus, cosinus, tangente, suite arithmétique, suite géométrique, intérêt simple, intérêt composé, Licence 1, Licence 2, CPTEC}}

% ==== Identité du document ====
\renewcommand{\DocTitre}{Mathématiques Générales (Licence 1 \& 2)}
\renewcommand{\DocKicker}{Institut Supérieur CPTEC \textbullet{} Support de cours}
\renewcommand{\DocSousTitre}{Des ensembles et de la logique aux suites numériques}
\renewcommand{\DocDescription}{Support de cours de l'étudiant : six chapitres avec exemples résolus}
\renewcommand{\BadgeUniq}{Licence 1 \& Licence 2}
\renewcommand{\DocVolumeH}{Année académique 2026--2027}

% ================= PAGE 1 : page de garde =================
\pagedegarde

% ================= PAGE 2 : page a l'effigie de CPTEC et de l'enseignant =================
\newcommand{\pagedeux}{%
\begin{titlepage}
\begin{tikzpicture}[remember picture, overlay]
% ================= bandeau superieur =================
\fill[robleu] (current page.north west)
  rectangle ([yshift=-10.8cm]current page.north east);
\fill[rogold] ([yshift=-10.8cm]current page.north west)
  rectangle ([yshift=-11.02cm]current page.north east);
\begin{scope}
  \draw[white, opacity=0.12, line width=1pt]
    ([xshift=14.2cm,yshift=-1.6cm]current page.north west) --
    ([xshift=16.6cm,yshift=-2.9cm]current page.north west) --
    ([xshift=19.3cm,yshift=-1.8cm]current page.north west) --
    ([xshift=18.4cm,yshift=-4.7cm]current page.north west) --
    ([xshift=15.6cm,yshift=-5.8cm]current page.north west) --
    ([xshift=13.8cm,yshift=-4.1cm]current page.north west) -- cycle;
  \draw[white, opacity=0.12, line width=1pt]
    ([xshift=14.2cm,yshift=-1.6cm]current page.north west) --
    ([xshift=18.4cm,yshift=-4.7cm]current page.north west);
  \draw[white, opacity=0.12, line width=1pt]
    ([xshift=16.6cm,yshift=-2.9cm]current page.north west) --
    ([xshift=15.6cm,yshift=-5.8cm]current page.north west);
\end{scope}
% ligne d'appel
\node[anchor=west, text=rogold!55!white, font=\small]
  at ([xshift=2.1cm,yshift=-1.85cm]current page.north west)
  {\textls[180]{\scshape\DocKicker}};
% titre principal sur deux lignes
\node[anchor=west, text=white, font=\fontsize{31}{38}\selectfont\bfseries]
  at ([xshift=2.0cm,yshift=-3.9cm]current page.north west) {MATHÉMATIQUES};
\node[anchor=west, text=white, font=\fontsize{31}{38}\selectfont\bfseries]
  at ([xshift=2.0cm,yshift=-5.55cm]current page.north west) {GÉNÉRALES};
\fill[rogold] ([xshift=2.05cm,yshift=-6.45cm]current page.north west)
  rectangle ++(2.8,-0.09);
% mention etudiant
\node[anchor=west, text=white, font=\Large\itshape]
  at ([xshift=2.1cm,yshift=-7.3cm]current page.north west) {Support de cours de l'étudiant};
\node[anchor=west, text=white!82!robleu, font=\small, text width=16.6cm, align=left]
  at ([xshift=2.1cm,yshift=-8.35cm]current page.north west)
  {Ensembles \textbullet{} Logique et algèbre de Boole \textbullet{} Équations et inéquations\\[1mm]
   Systèmes linéaires et programmation linéaire \textbullet{} Mesures des angles et trigonométrie\\[1mm]
   Suites numériques};
% ================= deux cartes : institut et enseignant =================
\begin{scope}[shift={(current page.north west)}]
% ---- carte 1 : l'institut ----
\fill[white, draw=rogris!45, line width=0.7pt, rounded corners=2.5mm]
  (1.8,-12.3) rectangle (10.2,-20.9);
\fill[robleu] (1.8,-12.3) rectangle (10.2,-13.0);
\node[text=white, font=\small\bfseries] at (6.0,-12.65)
  {\textls[140]{\scshape L'institut}};
\node[text=robleu, font=\large\bfseries, align=center, text width=7.6cm] at (6.0,-14.15)
  {CPTEC Business School\\[-0.5mm]{\footnotesize\mdseries Institut Supérieur}};
\node[text=roor, font=\footnotesize\bfseries] at (6.0,-15.15)
  {Agréé par l'État N\dgr{} 00.50/2024};
\node[text=rogris, font=\small, align=center, text width=7.4cm] at (6.0,-16.85)
  {Avedji, Lomé (Togo)\\[0.5mm]
   {\footnotesize Rue en face de l'Église Évangélique Presbytérienne\\ du Togo (EEPT), au niveau de SUNU Assurance}};
\node[text=rogris, font=\small, align=center, text width=7.4cm] at (6.0,-19.0)
  {\textbf{+228 70 78 78 18} \ \textbullet{} \ \textbf{96 51 05 50}\\[0.8mm]
   cptec228togo@gmail.com};
% ---- carte 2 : l'enseignant ----
\fill[white, draw=rogris!45, line width=0.7pt, rounded corners=2.5mm]
  (10.9,-12.3) rectangle (19.3,-20.9);
\fill[roteal] (10.9,-12.3) rectangle (19.3,-13.0);
\node[text=white, font=\small\bfseries] at (15.1,-12.65)
  {\textls[140]{\scshape L'enseignant}};
\node[text=robleu, font=\large\bfseries, align=center, text width=7.6cm] at (15.1,-14.3)
  {Ounimborbitibou DJABON};
\node[text=rogris, font=\small, align=center, text width=7.4cm] at (15.1,-16.15)
  {Titulaire de deux Masters :\\[0.4mm]
   Sciences Mathématiques (AIMS Ghana)\\[0.4mm]
   Physique Théorique (Université de Lomé)};
\node[text=rogris, font=\footnotesize\itshape, align=center, text width=7.4cm] at (15.1,-17.85)
  {Chercheur en Information et Communication Quantique};
\node[text=rogris, font=\small, align=center, text width=7.4cm] at (15.1,-19.35)
  {\textbf{+228 92 39 97 21}\\[0.8mm]
   djabon@aims.edu.gh\\[0.6mm]
   github.com/djabelo712};
\end{scope}
% ================= bandeau inferieur =================
\fill[robleu] ([yshift=-22.3cm]current page.north west)
  rectangle (current page.south east);
\fill[rogold] ([yshift=-22.3cm]current page.north west)
  rectangle ([yshift=-22.44cm]current page.north east);
\node[text=white, font=\fontsize{15}{18}\selectfont\bfseries]
  at ([yshift=-23.7cm]current page.north) {Niveau : Licence 1 \& Licence 2};
\node[text=white!85!robleu, font=\small]
  at ([yshift=-24.8cm]current page.north) {\textls[150]{\scshape\AnneeAcad}};
\end{tikzpicture}
\end{titlepage}}
\pagedeux

\pagenumbering{roman}

% numérotation des sections dans chaque chapitre (format du modele : 1, 2, 3 ...)
\renewcommand{\thesection}{\arabic{section}}

% ================= PAGE 3 : objectifs et table des matieres =================
{\color{robleu}\Large\bfseries Objectifs du module}\par
\vspace{1.5mm}

\textbf{Objectif général :} consolider les outils mathématiques fondamentaux (ensembles, logique, équations, systèmes linéaires et suites numériques) qui servent de socle à toutes les matières professionnelles : comptabilité, économie générale, techniques commerciales et informatique de gestion.

\vspace{2.2mm}
\textbf{Objectifs spécifiques :} à l'issue de cet enseignement, l'étudiant doit être capable de :
\begin{itemize}\small
  \item définir les ensembles, les représenter et appliquer les opérations entre ensembles ;
  \item raisonner avec la logique des propositions et simplifier avec l'algèbre de Boole ;
  \item résoudre dans $\R$ les équations et les inéquations du premier et du second degré ;
  \item résoudre les systèmes linéaires et formuler un problème de programmation linéaire ;
  \item mesurer les angles en degrés et en radians et utiliser les rapports trigonométriques ;
  \item calculer les termes et les sommes des suites arithmétiques et géométriques ;
  \item appliquer les suites aux intérêts simples et aux intérêts composés.
\end{itemize}

\vspace{4mm}
\begingroup
\makeatletter
\setcounter{tocdepth}{1}
\renewcommand*\l@chapter{\bfseries\@dottedtocline{0}{0em}{2.0em}}
{\noindent{\color{robleu}\Large\bfseries Table des matières}\par}\vspace{1.5mm}
{\small\@starttoc{toc}}
\makeatother
\endgroup

\clearpage
\pagenumbering{arabic}
